\documentclass[10pt,a4paper]{amsart}
\usepackage{amsmath,amssymb,amsthm,graphicx}
\usepackage{hyperref}
\numberwithin{equation}{section}

\newtheorem{Theorem}{Theorem}[section]
\newtheorem{Corollary}[Theorem]{Corollary}
\newtheorem{Lemma}[Theorem]{Lemma}
\newtheorem{Proposition}[Theorem]{Proposition}
{\theoremstyle{definition}
\newtheorem{Note}[Theorem]{Note}
\newtheorem{Definition}[Theorem]{Definition}
\newtheorem{Problem}[Theorem]{Problem}
}
\usepackage[a4paper,margin=22mm]{geometry}
\title{Commutativity of the q-Onsager elements xi\_1 and xi\_2 if and only if q\textasciicircum{}6 is not 1}
\author{Paul Terwilliger}
\date{doi:10.3842/SIGMA.2011.069; CC BY 4.0}
\begin{document}
\maketitle
\noindent\textit{Editorial note.} The counted claim is the author Proposition lem:RS2 at source lines 2199-2206, and its proof is the proof environment that immediately follows it at source lines 2208-2398; no other proof in the paper is used for it.\par\medskip
\noindent\textit{Editorial note.} The motivation prose at source lines 2099-2104 points to the alternative presentation Theorem prop:altpres, which is not part of this material; it is omitted because no step of the counted proof depends on it, not because any author argument was dropped.\par\medskip
\noindent\textit{Editorial note.} Definition def:qons carries an author citation in its optional argument; the bibliography fragment sources/sigma-2011-069-askey-wilson/qonsager-commutativity-bibliography.tex reproduces that cited entry verbatim from the source bibliography.\par\medskip
\noindent\textit{Allowed standard prerequisites (editorial).}\begin{itemize}
\item q-integer arithmetic in the standing field. The field F and the nonzero q in F with q\textasciicircum{}4 not 1 are as fixed at source lines 240-244, and the q-integer [n]\_q = (q\textasciicircum{}n - q\textasciicircum{}-n)/(q - q\textasciicircum{}-1) satisfies [3]\_q = 0 exactly when q\textasciicircum{}6 = 1.. Splits the two directions of the counted claim at source lines 2210-2211 and 2324-2326.
\item Free product of two cyclic groups of order three. For G = H * K with H and K cyclic of order three and generators h and k, every element of G has a unique reduced alternating normal form, and the elements of G form a basis of the group algebra F G.. Used at source lines 2343-2357 to define the module and at 2382-2392 to certify that the displayed element of F G is nonzero.
\item Associative expansion identities. The commutator identity of lem:gen(i)-(ii) is an expansion identity valid in any associative algebra, and a relation of the form X\textasciicircum{}3 Y - Y X\textasciicircum{}3 = c (X Y - Y X) makes the element X\textasciicircum{}3 - c X commute with Y, hence central in O.. Used at source lines 2330-2338 for the degenerate case q\textasciicircum{}6 = 1.
\end{itemize}\medskip
\begin{Definition}[\protect{\cite[Def\/inition~3.9]{qSerre}}]
\label{def:qons}
Let ${\mathcal O}={\mathcal O}_q$ denote the
 $\mathbb F$-algebra
 def\/ined by genera\-tors~$X$,~$Y$ and relations
 \begin{gather}
X^3Y-\lbrack 3\rbrack_q X^2YX+\lbrack 3\rbrack_q XYX^2-YX^3
= -\big(q^2-q^{-2}\big)^2(XY-YX),
\label{eq:td1}
\\
Y^3X-\lbrack 3\rbrack_q Y^2XY+\lbrack 3\rbrack_q YXY^2-XY^3
= -\big(q^2-q^{-2}\big)^2(YX-XY).
\label{eq:td2}
 \end{gather}
We call $\mathcal O$ the {\it $q$-Onsager algebra}.
\end{Definition}

\begin{Definition}
\label{lem:RS}
Let $\xi_1$, $\xi_2$ denote the following elements of
$\mathcal O$:
\begin{gather*}
\xi_1 =  XY-YX,
\\
 \xi_2 =
X^2Y^2-Y^2X^2+\big(q^2+q^{-2}\big)(YXYX-XYXY).
\end{gather*}
\end{Definition}

Referring to Def\/inition~\ref{lem:RS},
we are going to show that
$\xi_1$, $\xi_2$
commute if and only if~$q^6 \not=1$.
We will use the following results, which
apply to any $\mathbb F$-algebra.

\begin{Lemma}
\label{lem:gen}
%\label{lem:anyalg}
Let $x$, $y$ denote elements in any $\mathbb F$-algebra,
and consider the commutator
\begin{gather}
\big\lbrack
xy-yx, x^2y^2-y^2x^2 +\big(q^2+q^{-2}\big)(yxyx-xyxy)
\big\rbrack.
\label{lem:com12}
\end{gather}
\begin{enumerate}\itemsep=0pt
\item[$(i)$] The element
\eqref{lem:com12} is equal to
\begin{gather*}
 xyx^2y^2-x^2y^2xy+yxy^2x^2-y^2x^2yx
+
x^2y^3x-xy^3x^2+y^2x^3y-yx^3y^2
\\
 \qquad {} -
\big(q^2+q^{-2}\big)\big(
xyxy^2x
-xy^2xyx
+yxyx^2y
-yx^2yxy
\big).
\end{gather*}
\item[$(ii)$] The element
\eqref{lem:com12} times
$\lbrack 3\rbrack_q$ is equal to
\begin{gather}
\lbrack y,
\lbrack y,
\lbrack x,
\lbrack x,
\lbrack x,
y
\rbrack \,
\rbrack_q
\rbrack_{q^{-1}}
\rbrack_q
\rbrack_{q^{-1}}
+
\lbrack x,
\lbrack x,
\lbrack y,
\lbrack y,
\lbrack y,
x
\rbrack \,
\rbrack_q
\rbrack_{q^{-1}}
\rbrack_q
\rbrack_{q^{-1}},
\label{eq:double}
\end{gather}
where $\lbrack u,v\rbrack_\epsilon $ means $\epsilon uv-
\epsilon^{-1}vu$.
\end{enumerate}
\end{Lemma}

\begin{proof}
$(i)$ Expand
(\ref{lem:com12}) and simplify the result.
$(ii)$
Expand
(\ref{eq:double})
and compare it with the expression in $(i)$ above.
\end{proof}

\begin{Proposition}
\label{lem:RS2}
Referring to
Definition {\rm \ref{lem:RS}},
the elements $\xi_1$, $\xi_2$ commute
if and only if
$q^6 \not=1$.
\end{Proposition}

\begin{proof}
First assume
$q^6 \not=1$, so that
$\lbrack 3 \rbrack_q$ is nonzero.
Applying
Lemma
\ref{lem:gen}$(ii)$ to the elements $x=X$ and $y=Y$ in
the algebra $\mathcal O$, we f\/ind that
$\lbrack \xi_1, \xi_2\rbrack $ times $\lbrack 3\rbrack_q$
is equal to
\begin{gather}
\lbrack Y,
\lbrack Y,
\lbrack X,
\lbrack X,
\lbrack X,
Y
\rbrack \,
\rbrack_q
\rbrack_{q^{-1}}
\rbrack_q
\rbrack_{q^{-1}}
+
\lbrack X,
\lbrack X,
\lbrack Y,
\lbrack Y,
\lbrack Y,
X
\rbrack \,
\rbrack_q
\rbrack_{q^{-1}}
\rbrack_q
\rbrack_{q^{-1}}.
\label{eq:bigf}
\end{gather}
We show that the element (\ref{eq:bigf}) is zero.
Observe
\begin{gather*}
\lbrack X,
\lbrack X,
\lbrack X,
Y
\rbrack \,
\rbrack_q
\rbrack_{q^{-1}}
 =
X^3Y-\lbrack 3\rbrack_q X^2YX+\lbrack 3\rbrack_q XYX^2-YX^3
\\
\phantom{\lbrack X,
\lbrack X,
\lbrack X,
Y
\rbrack \,
\rbrack_q
\rbrack_{q^{-1}}}{} = -\big(q^2-q^{-2}\big)^2 \lbrack X,Y \rbrack
= \big(q^2-q^{-2}\big)^2 \lbrack Y,X \rbrack.
\end{gather*}
Similarly
\begin{gather*}
\lbrack Y,
\lbrack Y,
\lbrack Y,
X
\rbrack \,
\rbrack_q
\rbrack_{q^{-1}}
= \big(q^2-q^{-2}\big)^2 \lbrack X,Y \rbrack.
\end{gather*}
Therefore
\begin{gather*}
\lbrack Y,
\lbrack Y,
\lbrack X,
\lbrack X,
\lbrack X,
Y
\rbrack \,
\rbrack_q
\rbrack_{q^{-1}}
\rbrack_q
\rbrack_{q^{-1}}
 =
\big(q^2-q^{-2}\big)^2\lbrack Y,
\lbrack Y,
\lbrack Y,
X
\rbrack \,
\rbrack_q
\rbrack_{q^{-1}}
=
(q^2-q^{-2})^4\lbrack X,Y\rbrack.
\end{gather*}
Similarly
\begin{gather*}
\lbrack X,
\lbrack X,
\lbrack Y,
\lbrack Y,
\lbrack Y,
X
\rbrack \,
\rbrack_q
\rbrack_{q^{-1}}
\rbrack_q
\rbrack_{q^{-1}}
 =
\big(q^2-q^{-2}\big)^4\lbrack Y,X\rbrack .
\end{gather*}
By these comments
the element (\ref{eq:bigf})
is zero, and therefore
$\xi_1$, $\xi_2$ commute.


  Next assume
$q^6=1$, so that
$\lbrack 3\rbrack_q=0$ and
$(q^2-q^{-2})^2=-3$.
In this case the relations
(\ref{eq:td1}),
(\ref{eq:td2}) become
\begin{gather*}
X^3Y-YX^3=3(XY-YX),
\qquad
Y^3X-XY^3=3(YX-XY).
\end{gather*}
In the above line the relation on the left (resp. right)
asserts that $X^3-3X$ (resp.~$Y^3-3Y$)
commutes with $Y$ (resp.~$X$) and is therefore
central in~$\mathcal O$.
We show that $\xi_1$, $\xi_2$ do not commute by
displaying
an $\mathcal O$-module on which
$\lbrack \xi_1, \xi_2\rbrack$ is nonzero.
Def\/ine a group $G$ to be the free product
$H * K$ where each of~$H$,~$K$ is a cylic
group of order~3.
Let $h$ (resp.~$k$) denote a generator for~$H$ (resp.~$K$).
Let $\mathbb F G$ denote the group $\mathbb F$-algebra.
We give $\mathbb F G$ an $\mathcal O$-module structure.
To do this we specify the action of~$X$,~$Y$ on
the basis $G$ of $\mathbb F G$. For the identity $ 1\in G$
let $X.1=h$ and $Y.1=k$.
For $1 \not=g \in G$ write
$g=g_1g_2\cdots g_n$ such that
 for $1 \leq i \leq n$,
 $ g_i \in \lbrace h, h^{-1}\rbrace $
 for one parity of $i$ and
$g_i \in \lbrace k,k^{-1} \rbrace$ for the other parity of~$i$.
Let~$X$,~$Y$ act on~$g$ as follows:
\begin{center}
\begin{tabular}{c|c c c c}
Case &  $g_1=h$ & $g_1=h^{-1}$  & $g_1=k$ & $g_1=k^{-1}$ \\
\hline \hline
$X.g$ & $hg$ & $ 3h^{-1}g$ & $hg$ & $hg$ \\
$Y.g$ & $kg$ & $ kg$ & $kg$ & $3k^{-1}g$
    \end{tabular}
        \end{center}
We have now specif\/ied the actions of $X$, $Y$
on~$G$.
These actions
give $\mathbb F G$ an
 $\mathcal O$-module
structure on
which $X^3=3X$ and $Y^3=3Y$.
For the $\mathcal O$-module $\mathbb F G$ we now
apply
$\lbrack \xi_1, \xi_2\rbrack$ to the vector $1$.
By
Lemma~\ref{lem:gen}$(i)$ and the construction,
the element
$\lbrack \xi_1, \xi_2\rbrack$ sends  $1$
to
\begin{gather*}
 hkh^{-1}k^{-1}-h^{-1}k^{-1}hk+
khk^{-1}h^{-1}-k^{-1}h^{-1}kh
+
  3 h^{-1}kh-3hkh^{-1}
+ 3 k^{-1}hk-3khk^{-1}
\\
\qquad {} +
hkhk^{-1}h-hk^{-1}hkh+
khkh^{-1}k-kh^{-1}khk.
\end{gather*}
The above element of $\mathbb F G$
is nonzero. Therefore
$\lbrack \xi_1, \xi_2\rbrack \not=0$
so
$\xi_1$, $\xi_2$ do not commute.
\end{proof}

\begin{thebibliography}{99}
\bibitem{qSerre}
Terwilliger P.,
Two relations that generalize the $q$-Serre relations and the  Dolan--Grady relations,
in  \href{http://dx.doi.org/10.1142/9789812810199_0013}{Physics and Combinatorics 1999} (Nagoya), World Sci. Publ., River Edge, NJ, 2001, 377--398,
\href{http://arxiv.org/abs/math.QA/0307016}{math.QA/0307016}.
\end{thebibliography}
\end{document}
